\section{Result Records, Comparability Rules and Validity}\label{sec:scope}
The atomic unit of our repository is one \emph{evaluated result record}, rather than a whole paper or a single headline score. A result record binds a source version and locator to seven structured fields: Source, Target (and threat polarity Types~I--IV), Receiver projection, Intervention boundary, Attacker training/calibration ($D_0$: clean-condition fit; $D_2$: protected-condition fit), Workload/holdout, and Outcomes (inference metric, task/component utility, and cost object; Table~\ref{tab:result-record}). A study family contributes multiple records whenever it evaluates distinct execution modes (e.g., stored capture vs.\ live loopback in Dangol \cite{dangol2026workflow}), distinct mechanisms (e.g., $K$-prefix KV simulation vs.\ GPTCache semantic de-identification in Early Bird \cite{earlybird2025}), or separately scoped performance benchmarks (e.g., SafeKV \cite{chu2025safekv} and OTRO \cite{lee2026otro}).
\begin{table}[t]\centering\footnotesize
\caption{A single result record. Only fields relevant to the requested comparison must align; unavailable required fields remain unresolved.}\label{tab:result-record}
\begin{tabularx}{\columnwidth}{@{}p{.22\columnwidth}Y@{}}\toprule
Record field & Required meaning\\\midrule
Source & Family/result ID, version, source locator, reading scope and hash.\\
Target & Hidden property, inference goal and protection relation: confidentiality, attribution, human/agent distinction or authorized diagnosis.\\
Receiver & Actual permission and released representation, including session/document linkage and any auxiliary inputs.\\
Intervention & Where protection ends; whether edits affect stored traces, transmission or executed actions.\\
Attacker & Fit/calibration condition, test population, training access and adaptation budget.\\
Workload & Task, system, topology/hardware and grouping/holdout.\\
Outcomes & Inference metric / denominator; objective or proxy utility; applicable cost object; uncertainty and field states.\\\bottomrule
\end{tabularx}\end{table}


\paragraph{Indexing schema and comparison contracts}
The mnemonic \scope{} (Secret/target, Channel, Observer, Pipeline, Exploitation goal) names the target-and-channel columns of our index, paired with evaluation provenance. Generating layers span hardware (P1), serving (P2), transport (P3), application/API (P4), orchestration (P5), tools/MCP (P6), environment (P7) and multi-agent communication (P8). We apply three comparison contracts (Table~\ref{tab:comparison-contracts}): a comparison is \emph{compatible} when required fields align; \emph{incompatible} when a positive mechanism or boundary conflict is established (such as substituting token-payload ratios for wire bytes or offline timestamp edits for live action delays); and \emph{unresolved} when a required binding is unreported (U), unavailable (A) or unverified (?). Table~\ref{tab:result-lineage} illustrates a complete record-to-judgment lineage.
\begin{table}[t]\centering\footnotesize
\caption{Comparison rules. A known mismatch and a missing binding have different consequences.}\label{tab:comparison-contracts}
\begin{tabularx}{\columnwidth}{@{}p{.28\columnwidth}Y@{}}\toprule
Requirement & Condition and resulting license\\\midrule
Descriptive comparison & Align target, receiver permission, metric meaning and cost object. Different systems may be described together; token payload and wire bytes cannot be ranked as the same cost.\\
Paired effect & Compare the same defended test units and outcome. A known test-unit mismatch makes that subtraction incompatible; an unverified fit binding leaves the effect unresolved.\\
Deployment transfer & Establish a bridge for the changed receiver, intervention, workload, topology or hardware. A demonstrated violation blocks the proposed transfer; an untested bridge leaves it unresolved. Permit only the range and endpoints actually covered.\\\bottomrule
\end{tabularx}\end{table}

\begin{table}[t]\centering\footnotesize
\caption{Filled record and its licensed judgment. Values are source-reported; no cross-mode effect is computed.}\label{tab:result-lineage}
\begin{tabularx}{\columnwidth}{@{}p{.24\columnwidth}Y@{}}\toprule
Join key / field & Entry\\\midrule
Family / result & Dangol workflow / \texttt{DANGOL-livebinding}.\\
Version / locator & arXiv:2606.07150v3, \S9.7, live-protocol paragraph after Table~5 \cite{dangol2026workflow}.\\
Target / view & Workflow classification from live Oblivious HTTP (OHTTP)/Agent-to-Agent (A2A) traffic; not GUI-state observation.\\
Intervention & Live fixed-rate cover with fresh per-session identifiers, distinct from transformed stored captures.\\
Workload / testbed & Workflow classes with client, relay, gateway and agent co-located on one loopback host.\\
Outcome / cost & Accuracy 0.30; about $13\times$ bytes and $11\times$ latency in that condition.\\
Unresolved fields & Objective terminal task utility \NRR{}; precise paired live-refit comparison \NRV{}.\\
Licensed judgment & These costs describe this live protocol condition. They neither price dummy GUI actions nor prove task preservation. A matched objective endpoint is needed for the latter.\\\bottomrule
\end{tabularx}\end{table}



\paragraph{Executable comparison rules}
The record schema and comparison contracts precede the audit results. Table~\ref{tab:audit-rules} gives all nine operational rules and a shared evidence-status alphabet; Table~\ref{tab:family-audit} applies them to every mitigation family. Twenty-four constructed cases and seven mechanism checks exercise rule application in the artifact, while the source-linked cases in \S\ref{sec:discussion} explain the substantive comparison judgments.
